% Statement on the use of generative AI, with the chapter-by-chapter account
% folded in beneath it. Back matter: follows the originality index and
% immediately precedes the pooled references. Treated as a \chapter*, like
% "Originality of the results" -- not as an appendix.
%
% The internal heading uses \oisec (from originality_index/setup) so that the
% two back-matter chapters share a house style.

\chapter*{Statement on the use of generative AI}
\phantomsection
\label{sup:ai-use}
\addcontentsline{toc}{chapter}{Statement on the use of generative AI}
\markboth{STATEMENT ON THE USE OF GENERATIVE AI}{STATEMENT ON THE USE OF GENERATIVE AI}

I have made extensive use of generative AI and agentic coding tools in producing this thesis, and it is better to be forthcoming and candid about that than to leave it implicit. We live in a rapidly changing digital landscape that is increasingly challenging preconceptions, and it falls to our generation to help define the standard for how AI is taught and integrated into higher education. To that end, I give a summary below of how I used AI in the writing of this thesis, with a chapter-by-chapter account following it.

Where AI produced mathematics, prose or code, I have checked it and I take full responsibility for the content of this thesis.

\oisec{Account of generative AI use}
\phantomsection
\label{sup:ai-account}

The following is a short description of how generative AI was used to assist the production of each chapter.

\paragraph{Chapter 1.}
There was a full introduction in the pre-AI draft of 2024. After the other nine chapters had all been fully redrafted, Claude was used quite simply to modify the existing introduction to reflect the changes in presentation of the new chapters.

\paragraph{Chapter 2.}
Chapter 2 was partially taken from the pre-AI draft, and partially added to with agent tools. All of the MOC content and much of the branching process definition was in the original. But the Markov Chains section was largely generated.

\paragraph{Chapter 3.}
The Koenigs result was initially determined by ChatGPT in January 2026. At that time, LLMs were not quite good enough to write sufficient quality prose for the thesis, so I wrote up the whole content of this chapter myself. Later, in July 2026, AI was more capable and I used Claude to rework the existing content into a better flow, and also used Codex to build the fractal generation web tool.

\paragraph{Chapter 4.}
The original chapter was re-worked in a style pass. The same content from the pre-AI draft was presented more coherently with a small addition of the $D_t$ variable.

\paragraph{Chapter 5.}
The old BDC extension chapter was divided into what are now Chapters 5 and 6. Chapter 5 is largely the same story as Chapter 4, in that the same content was re-structured and presented better.

\paragraph{Chapter 6.}
This chapter extends the old ``BDC extension'' chapter, whose content was clearly unresolved. Among the key additions are the BMVR-style bursting equations, and the $L$ parameter flooding analysis. Practically all of this additional content has been generated autonomously with LLMs. And for this chapter more than any other, the same can be said of the \LaTeX\ write-up.

\paragraph{Chapter 7.}
The crux of this chapter was present in the pre-AI draft. I had previously shown that the two-type BDC system can be solved with the specified method. I had merely neglected to actually write down the final solution explicitly. There was some complication to that task, but Claude did it effortlessly, and restructured my prose along the way.

\paragraph{Chapter 8.}
The existing content was restructured and styled with LLM passes.

\paragraph{Chapter 9.}
Content was trimmed and rearranged into a variable time model collection.

\paragraph{Chapter 10.}
The quadrant argument was written by me, apart from its figure; everything else in the chapter was generated automatically.

\medskip
I have used Claude Code (Pro), Grok Build (SuperGrok), Codex (Plus),  and Zcode (GLM Coding Lite). Claude and Codex are the most capable models; Grok and Zcode are cheaper and faster, but still reliable alternatives for more mechanical tasks. The originality index was fully researched, verified, and written up by Grok and Claude.

The multiplicative dissipation cellular automaton (CA) in Chapter 9 was agentically coded end-to-end. But other substantial programs were written by me. Among these is the agent-based simulation shown in Chapter 1 (made with p5.js).

There is also the burst suppression simulator and the Pimentel CA. Both of these were fully implemented in C++ using SFML (Simple and Fast Multimedia Library), then later ported to JavaScript and given a new user interface.

Overall, I can say that AI has been remarkably helpful in producing a final result greatly superior to the pre-AI draft. And that generally, LLMs and coding agents have been really enjoyable to use and have recently become surprisingly capable.

If you would like to know more about how I have been using AI, please feel free to get in touch, and I would be happy to talk about it with you.
